\section{Failure analysis}
\label{sec:failure}

% Figure 8 -- observed failure categories by condition.
% Data: figures/data/fig8_failure_mix.dat, from PHASE3/results/per_model_config.csv (primary models).
\begin{figure}[t]
\centering
\begin{tikzpicture}
\begin{axis}[
  paperaxis,
  width=0.95\linewidth, height=5.6cm,
  ybar stacked, bar width=7.5pt,
  ymin=0, ymax=100,
  ylabel={\% of runs},
  symbolic x coords={B,G,T,R,RT,LA,LG,EB,EBG,EBG-all,EBG-unsafe,EBG-noR3},
  xtick=data,
  x tick label style={font=\scriptsize, rotate=32, anchor=east},
  ymajorgrids, ytick={0,20,40,60,80,100},
  legend style={at={(0.5,-0.42)}, anchor=north, legend columns=5, draw=cNeutral},
  legend cell align=left,
]
\addplot[draw=cEBG, fill=cEBG, fill opacity=0.85]
  table[col sep=tab, x=config, y=correct] {figures/data/fig8_failure_mix.dat};
\addplot[draw=cAdv!70!black, fill=cAdv, fill opacity=0.9]
  table[col sep=tab, x=config, y=false_conflict] {figures/data/fig8_failure_mix.dat};
\addplot[draw=cNone, fill=cNone, fill opacity=0.55]
  table[col sep=tab, x=config, y=over_abstain] {figures/data/fig8_failure_mix.dat};
\addplot[draw=cHarm, fill=cHarm, fill opacity=0.9, postaction={pattern=north east lines, pattern color=white}]
  table[col sep=tab, x=config, y=harmful] {figures/data/fig8_failure_mix.dat};
\addplot[draw=cLLM!70!black, fill=cLLM, fill opacity=0.9]
  table[col sep=tab, x=config, y=format_err] {figures/data/fig8_failure_mix.dat};
\legend{correct, false conflict, unnecessary abstention, harmful answer, format failure}
\end{axis}
\end{tikzpicture}

\vspace{4mm}
\caption[Observed failure categories by condition]{%
\textbf{Where the failures go.} Outcome rates from the frozen scorer, pooled over the four primary models and four
decodings ($n=4288$ runs per condition). Only \emph{observed} categories are shown; each is a rate over all runs,
and because the scorer's categories are not mutually exclusive the five bands do not sum to 100\% --- the
remainder (step-limit exhaustion, stale values counted inside ``harmful'', and unclassified residue) is omitted
rather than invented. Three readings matter. Metadata labels (\cond{G}) leave the agent \emph{worse} than no labels
(\cond{B}): more false conflicts and more harm. Enforcing record-derived labels (\cond{R}, \cond{RT}) produces the
largest harmful bands in the study. And the residual failure of the proposed condition \cond{EBG} is
over-conservatism --- 20.2\% false conflicts and 11.5\% unnecessary abstentions --- against 3.1\% harm; the
mechanism moved errors where it was designed to move them, which is a claim about direction, not about whether
that is the right place for them.\\
\emph{Accessible description:} twelve stacked bars; the correct band is largest for \cond{EBG-all} and
\cond{EBG}, while the hatched harmful band is widest for \cond{RT}, \cond{LG} and \cond{R} and almost invisible for
\cond{EB} and \cond{EBG}.\\
\emph{Source:} \srcfile{PHASE3/results/per_model_config.csv}; plotted from
\srcfile{figures/data/fig8_failure_mix.dat}.}
\label{fig:failures}
\end{figure}


\subsection{Where the failures go}

\cond{EBG} answers 61.3\% of pooled primary runs correctly, so roughly two runs in five still fail. Of those,
20.2\% of all runs report a false conflict, 11.5\% abstain unnecessarily, 2.8\% answer with the superseded value,
2.8\% produce no parsable answer, 1.0\% exhaust the step budget and 0.3\% are arbitrary picks; no run produced an
unsupported claim and none was missing.

The residual failure is therefore \textbf{over-conservatism, not harm}: false conflicts and unnecessary
abstentions together account for 31.7\% of runs against 3.1\% harmful. The mechanism was designed to move errors
in exactly this direction, and it did. Whether that is the right place for them is a deployment question this
benchmark cannot settle, because it assigns no relative cost to a wrong confident answer and a missed correct one.

Part of the false-conflict residue is the arbiter's own documented limitations, pinned by tests before the freeze
and deliberately not fixed during the campaign: a negator more than four tokens before the cue is missed,
``correcting'' is not treated as a correction cue, and one negated-both-stand pattern resolves to the rival instead
of a contradiction. We cannot say how much of the 20.2\% these account for without a transcript study
(\Cref{sec:limitations}).

\subsection{The two mechanism-level failures the design targets}

False resolution --- asserting an ordering the evidence does not support --- falls from 12.8\% of conflict groups
under plain memory, 19.0\% under \cond{RT} and 37.0\% under \cond{LG} to \textbf{1.9\%} under \cond{EBG}. Conflict
preservation rises from 65.2\%, 66.8\% and 42.4\% to \textbf{88.5\%}. These two measures, rather than task
success, are where the mechanism's design is visible most directly.

\subsection{Model families fail differently in kind}

\Cref{fig:failures} is pooled; the per-model picture divides along family lines. Under \cond{EBG} the three Qwen
models abstain unnecessarily on 0.0--3.6\% of runs and fail to format an answer on at most 0.1\%; Llama-3.1-8B
abstains on 44.0\% and fails to format on 11.1\%, while being the \emph{least} harmful model in the study under
every condition. It rarely commits to a wrong value because it rarely commits. This is the sharpest limit on
generalisation in our results: the harm-reduction claim is demonstrated on models that actually produce harmful
answers. Qwen2.5-14B's 13.8\% format-error rate is constant across all twelve conditions and is a harness
formatting artefact, not a mechanism effect.

\subsection{Execution failures}

There were none in the analysis sense: all 18 cells reached 3216 rows with 0 missing pairs, 0 duplicates and 0
harness-error rows across all 57{,}888 runs, so no run was scored as missing and the pre-registered missing-run
rule never applied. Two cells were interrupted by the infrastructure event described in
\Cref{sec:limitations} and resumed without losing a completed run.

\subsection{What this analysis does not cover}

The eval3 task texts, gold labels and per-run transcripts remain sealed until publication, so every statement here
is at the level of the frozen scorer's outcome categories and their rates. We did not inspect individual failing
runs. In particular we cannot say which of \cond{EBG}'s false conflicts the arbiter could have resolved from the
text and which are items where reporting a conflict is correct, nor can we confirm that the development failure
mechanisms are the ones that produced the eval3 residue. Both require a transcript study after the data release.
